\documentclass[11pt]{article}
\usepackage[utf8]{inputenc}
\usepackage[T1]{fontenc}
\usepackage[ngerman]{babel}
\usepackage[a4paper, margin=2.5cm]{geometry}
\usepackage{graphicx} % Required for inserting images
\usepackage{parskip}
\pagestyle{empty}
\setlength{\parindent}{0pt}

\title{Bewerbung}
\author{Felix Dietrich} 
\date{October 2026}

\begin{document}

\textbf{Bewerbung als Doktorand im Bereich Künstliche Intelligenz in der Fahrzeugentwicklung}

\bigskip

Sehr geehrte Damen und Herren,

als zukünftiger Masterabsolvent der Informatik an der Universität Tübingen bin ich auf der Suche nach einer Doktorandenstelle, in der ich meine Kenntnisse aus dem Studium wissenschaftlich vertiefen und zugleich auf reale Fragestellungen in der Industrie anwenden kann. Die ausgeschriebene Position im Bereich Künstliche Intelligenz in der Fahrzeugentwicklung bietet hierfür eine besonders attraktive Möglichkeit, da sie wissenschaftliche Forschung mit der Entwicklung datengetriebener Lösungen für konkrete Fragestellungen der Fahrzeugentwicklung verbindet.

In meinem Bachelorstudium der Wirtschaftsinformatik an der Hochschule Reutlingen beschäftigte ich mich intensiv mit der Schnittstelle zwischen betrieblichen Prozessen und Informationstechnologie und sammelte frühzeitig praktische Erfahrung in der Softwareentwicklung und Datenverarbeitung. Während meines Praxissemesters bei der Mercedes-Benz AG arbeitete ich direkt mit Ingenieuren der Fahrzeugentwicklung zusammen und entwickelte eine VBA-basierte Softwarelösung, die wiederkehrende Arbeitsschritte automatisierte und die Qualitätssicherung im Entwicklungsprozess unterstützte. Dadurch trug ich dazu bei, den Entwicklungsprozess von Fahrzeugen zu beschleunigen und Ingenieuren mehr Zeit für fachliche Bewertungen und Entscheidungen zu geben. Für die Robert Bosch GmbH entwickelte ich während meines Bachelorstudiums eine Engineering-Data-Pipeline für CAN-Bus-Daten mit Python weiter. Im Mittelpunkt stand dabei die automatisierte Fehleranalyse von Daten aus Erprobungsfahrten. Die Lösung unterstützte die Auswertung großer Datenmengen und half, Auffälligkeiten systematisch zu identifizieren. Damit leistete sie einen Beitrag dazu, die Entwicklung und Optimierung von Bauteilen datenbasiert zu unterstützen und Erkenntnisse aus realen Erprobungsfahrten schneller in den Entwicklungsprozess einzuspeisen.

Um meine Kenntnisse in modernen Informationstechnologien zu vertiefen, nahm ich anschließend den Masterstudiengang Informatik an der Universität Tübingen auf. Dort beschäftigte ich mich unter anderem mit maschinellem Lernen, neuronalen Netzen und Zeitreihenanalyse. Diese Kenntnisse setze ich derzeit in meiner Masterarbeit praktisch ein, in der ich eine cloud-native Anwendung auf AWS mit einer Machine-Learning- und Optimierungskomponente entwickle. Im Mittelpunkt steht ein Online-Optimierungsverfahren auf Basis eines Gauß-Prozess-Modells. Das Modell wird mit scikit-learn auf Amazon SageMaker anhand fortlaufend erhobener Nutzungsdaten trainiert und verbessert Parameter einer Anwendung im laufenden Betrieb iterativ. Wenn ich die Arbeit ende Oktober finalisiert habe kann ich sie ihnen gerne zukommen lassen, damit Sie sich von meiner Sorgfalt und Leidenschaft für Projekte die Praxis mit wissentschaftlichem Arbeiten verknüpfen überzeugen könnnen.

Meine bisherigen Stationen bei Porsche haben meinen Wunsch bestärkt, meinen weiteren beruflichen Weg langfristig bei Porsche zu gestalten. Als Teil der Cellforce Group konnte ich bereits ein fachlich anspruchsvolles und zugleich sehr wertschätzendes Arbeitsumfeld kennenlernen. Dieser positive Eindruck bestätigt sich in meiner aktuellen Tätigkeit als Werkstudent bei der Porsche AG in Zuffenhausen. Die Promotion bietet mir die Möglichkeit, meine Kenntnisse in Künstlicher Intelligenz, Zeitreihenanalyse und Datenanalyse mit meinem Verständnis für betriebswirtschaftliche Abläufe zu verbinden und daraus einen konkreten Beitrag zur Weiterentwicklung einer zukunftsorientierten und KI-gestützten Fahrzeugentwicklung zu leisten.

Über die Gelegenheit, Sie in einem persönlichen Gespräch von meiner Motivation und meiner fachlichen Eignung zu überzeugen, freue ich mich sehr.

\bigskip

Mit freundlichen Grüßen

\medskip

Felix Dietrich

\end{document}
